\documentclass[10pt]{article}
\usepackage[T1]{fontenc}
\usepackage{lmodern,amsmath,amssymb,amsthm,microtype}
\usepackage[margin=.82in]{geometry}
\usepackage[colorlinks=true,urlcolor=blue]{hyperref}
\newtheorem{theorem}{Theorem}
\title{Column Differences and Alternating Antidiagonals of A260056}
\author{Qichao Wang\\{\small Hebei University of Technology}\\[.3em]Daoyu Dong\\{\small University of Electronic Science and Technology of China}}
\date{October 26, 2026}
\begin{document}\maketitle
\begin{abstract}
We prove both conjectures recorded in OEIS A260056. Every column is
an explicit polynomial in the row index, of degree exactly one more
than the column index. This proves the conjectured finite-difference
recurrence and its minimal order. Alternating antidiagonal sums have
a rational generating function that identifies them with the stated
shift of A004524. All assertions hold at every row and column.
\end{abstract}
Let
\[
 T(n,k)=[x^k]\sum_{j=0}^n(1+x+x^2)^j\quad(n,k\ge0),
\]
with coefficients outside a row taken as zero. This is the polynomial
definition of A260056~\cite{target}. The entry conjectures that column
$k$ has vanishing differences of order $k+2$, and that its alternating
antidiagonal sums equal A004524 shifted by three.
\begin{theorem}
For every $n,k\ge0$,
\begin{equation}\label{eq:column}
 T(n,k)=\sum_{i=\lceil k/2\rceil}^{k}
   \binom{i}{k-i}\binom{n+1}{i+1}.
\end{equation}
Consequently
\[
 \sum_{m=0}^{k+2}(-1)^m\binom{k+2}{m}T(n+m,k)=0.
\]
As a polynomial in $n$, column $k$ has degree exactly $k+1$;
no lower positive difference order vanishes identically.
\end{theorem}
\begin{proof}
Expand $(1+x+x^2)^j=\sum_{i=0}^j\binom ji x^i(1+x)^i$.
Its coefficient of $x^k$ is
$\sum_i\binom ji\binom i{k-i}$, where only
$\lceil k/2\rceil\le i\le k$ contribute. Sum over $j$ and apply
$\sum_{j=0}^n\binom ji=\binom{n+1}{i+1}$ to obtain
\eqref{eq:column}, including when $i>n$.
The summand $i=k$ has degree $k+1$ and leading coefficient
$1/(k+1)!$; every other summand has smaller degree.
Each forward difference lowers the degree of a nonconstant
polynomial by exactly one. Thus the $(k+2)$nd difference is zero
and the $(k+1)$st difference is the nonzero constant one.
The displayed signed form differs from the forward-difference
convention only by the global sign $(-1)^{k+2}$.
\end{proof}
\begin{theorem}
Define $c(m)=\sum_{k=0}^m(-1)^kT(m-k,k)$. Then for every $m\ge0$,
\[
 c(m)=\mathrm{A004524}(m+3),\qquad
 c(4q+r)=2q+\begin{cases}1,&r=0,\\2,&r=1,2,3,\end{cases}
 \quad 0\le r<4.
\]
\end{theorem}
\begin{proof}
The row definition gives the bivariate generating function
\[
 F(t,x)=\sum_{n,k\ge0}T(n,k)t^nx^k
 =\frac1{(1-t)(1-t(1+x+x^2))}.
\]
Substitution $(t,x)=(z,-z)$ is legitimate as a formal series
and weights each term by $(-1)^kz^{n+k}$. Hence
\[
 \sum_{m\ge0}c(m)z^m=F(z,-z)
 =\frac1{(1-z)^2(1+z^2)}.
\]
A004524~\cite{comparison} has generating function
$z^3/((1-z)^2(1+z^2))$ and initial coefficients zero through
index two, proving the stated shift. Alternatively, the displayed
four residue-class formulas have the same generating function,
which verifies the identification without relying on a numerical
comparison of terms.
\end{proof}
\paragraph{Validation and scope.}
The standard-library verifier independently constructs rows using
trinomial convolution, checks the closed column formula through row
100 and column 40, audits the finite-difference identities, checks
every antidiagonal through 100, and compares the complete published
100-row flattened table. The arguments prove the two conjectures
for all indices; finite checks validate the implementation.
\begin{thebibliography}{2}
\bibitem{target} W. Schulte, OEIS A260056,
\url{https://oeis.org/A260056}, accessed October 26, 2026.
\bibitem{comparison} The OEIS Foundation, A004524,
\url{https://oeis.org/A004524}, accessed October 26, 2026.
\end{thebibliography}\end{document}
